\chapter{Kustomize: Manifests as a Set}
\label{ch:kustomize}

\section{The problem}

Part~V produced a dozen YAML files. Applying them one by one invites
forgetting one; and the moment a second environment exists (``prod, but
with 2 replicas and real secrets''), copy-pasting manifests per
environment guarantees drift. Two tools dominate this space; this
project uses the one built into kubectl.

\section{The kustomization, annotated}

\code{kubectl apply -k deploy/k8s/base} reads one file:

\begin{lstlisting}[language=yaml]
# deploy/k8s/base/kustomization.yaml
namespace: ts                       (*@\co{1}@*)
resources:                          (*@\co{2}@*)
  - postgres-course.yaml
  - postgres-auth.yaml
  - mongodb.yaml
  - minio.yaml
  # kafka.yaml requires the Strimzi operator (one-time admin step:
  # ../install-strimzi.sh) -- without its CRDs, apply fails on kind Kafka.
  - kafka.yaml
  - kafka-topics.yaml
  - maildev.yaml
  - config-server.yaml
  - discovery-server.yaml
  - course-service.yaml
  - auth-service.yaml
  - dictionary-service.yaml
  - notification-service.yaml
  - api-gateway.yaml
commonLabels:                       (*@\co{3}@*)
  app.kubernetes.io/part-of: teacher-supporter
\end{lstlisting}

\begin{description}
  \item[\co{1}] Stamped onto every resource --- no manifest needs its
  own \code{namespace:} line, and retargeting the whole app to another
  namespace is a one-line change.
  \item[\co{2}] The definitive list. Two deliberate absences teach more
  than the presences: \code{namespace.yaml} is out because it is
  cluster-scoped and Jenkins' namespaced Role could not apply it
  (Chapter~22); Secrets are out for Chapter~16's reasons. A third
  object with the same admin-only status is not a file at all: the
  Strimzi operator and its CRDs, installed once by
  \code{install-strimzi.sh} (Chapter~25) --- \code{kafka.yaml} is listed
  here but cannot apply until they exist. The listed order is
  documentation only --- the API server accepts everything and
  reconciliation sorts out reality.
  \item[\co{3}] A label stamped on all objects, making the whole
  application one selector: \code{kubectl get all -l
  app.kubernetes.io/part-of=teacher-supporter} --- or, in a disaster
  drill, one delete.
\end{description}

\section{Bases and overlays}

Kustomize's real feature arrives with the second environment. A
\emph{base} holds the truth; \emph{overlays} patch it:

\begin{lstlisting}[language=yaml]
# deploy/k8s/overlays/prod/kustomization.yaml
resources:
  - ../../base
patches:
  - target: { kind: Deployment, name: course-service }
    patch: |-
      - op: replace
        path: /spec/replicas
        value: 2
images:
  - name: localhost:5000/course-service
    newTag: v1.4.2
\end{lstlisting}

No templates, no variables --- overlays are \emph{merges of plain YAML}.
Every layer is valid Kubernetes YAML you can read, diff, and pipe
through \code{kubectl kustomize} to inspect the final output before
anything touches the cluster. The directory layout is already
overlay-shaped (\code{base/}), so the day a prod overlay is needed, it
is an \code{mkdir}, not a migration.

\begin{decision}[Kustomize, not Helm]
Helm is the other answer: manifests as Go templates
(\code{replicas: \{\{ .Values.replicaCount \}\}}) with per-environment
value files, packaged as versioned, shareable \emph{charts}. Templating
is strictly more powerful --- loops, conditionals, one chart deploying
anyone's app --- and that is the point of divergence: Helm shines when
you \emph{distribute} software to strangers' clusters; for deploying
your own app, templates make every manifest unreadable until rendered.
Kustomize keeps plain YAML plain, ships inside kubectl, and covers
replicas/images/labels patching --- the whole need here. You will still
\emph{consume} Helm charts for third-party software (Strimzi,
kube-prometheus-stack); writing one can wait until you have users.
\end{decision}

\begin{pitfall}[spec.selector: field is immutable]
Symptom: the pipeline's \code{kubectl apply -k} fails on one
Deployment with this error, although the same file applies fine by
hand. Cause: \code{commonLabels} stamps its label not only onto
metadata but into every Deployment's \emph{selector} --- and selectors
are immutable once created. An object first created \emph{outside}
Kustomize (\code{kubectl apply -f maildev.yaml} during milestone~6's
setup, selector \code{\{app: maildev\}}) can never be updated by
Kustomize afterwards (selector
\code{\{app: maildev, part-of: teacher-supporter\}}). Fix: delete the
hand-made object and let the pipeline recreate it. Rule: an object
managed by a kustomization is applied \emph{only} through it --- for a
one-off test, render with \code{kubectl kustomize} and pipe that,
never the raw file. (Modern Kustomize's \code{labels:} with
\code{includeSelectors: false} avoids selector injection --- but
switching now would break every selector already created the old way.)
\end{pitfall}

\begin{pitfall}[it worked yesterday: the prune illusion]
Symptom: you delete a manifest from \code{resources:}, re-apply, and
the object is still running --- for months, until it conflicts with
something. Cause: apply-the-set is still apply --- Chapter~14's rule
that apply never deletes holds for Kustomize too. The removed Ingress
of Chapter~4 needed a manual \code{kubectl delete}. The systematic
fixes: \code{kubectl apply -k --prune} (still gated behind flags for
good reason --- a mislabeled selector can mass-delete), or GitOps
controllers (Chapter~27) whose whole job is making git and cluster
converge, deletions included.
\end{pitfall}

\begin{handson}[render before you trust]
\begin{lstlisting}[language=bash]
kubectl kustomize deploy/k8s/base | less     # the final YAML, merged
kubectl kustomize deploy/k8s/base | grep -c 'kind:'   # object count
kubectl apply -k deploy/k8s/base --dry-run=server     # validate, no changes
\end{lstlisting}
\code{kubectl kustomize} is pure text transformation --- reading its
output is how you review what an overlay actually changes, and
\code{--dry-run=server} is the closest thing to a manifest unit test.
\end{handson}

\begin{quiz}
\begin{enumerate}
  \item What two problems does a kustomization solve over
  \code{kubectl apply -f} per file, even with zero overlays?
  \item Why are the namespace and the Secrets deliberately missing from
  \code{resources:}? (Two different reasons.)
  \item Helm vs.\ Kustomize: state each tool's home turf in one
  sentence, and name the third-party software in this project's future
  that will arrive as a chart.
  \item A teammate removes \code{mongodb.yaml} from the list to
  ``decommission'' it. What actually happens, and what are the two
  correct procedures?
\end{enumerate}
\end{quiz}
